\section{System Overview}
\label{sec:system}

The pipeline separates perception, identity, and presentation geometry:
\begin{equation}
\begin{aligned}
\text{video} &\rightarrow \text{boxes + pose keypoints}\\
&\rightarrow \text{local tracklets}\\
&\rightarrow \text{offline identity paths}\\
&\rightarrow \text{stable rider views}.
\end{aligned}
\label{eq:system-pipeline}
\end{equation}

The detector is a custom-trained one-class YOLO11m-pose model.  It was initialized
from public pretrained weights and trained for 700 epochs at 640-pixel input
resolution on manually reviewed boxes and two windsurf-specific pose keypoints: the
boom--mast junction and mast tip.  Each box is also converted into a compact sail-color
descriptor after border-connected background suppression.  Independently, background features
are estimated between adjacent frames while masking detected surfers; the resulting
camera transforms are passed to both tracking stages.

The first association stage is causal but deliberately conservative.  It accepts
only locally unambiguous matches and exposes uncertainty as separate tracklets.  The
second stage sees the complete recording, creates a bounded graph of feasible
tracklet continuations, and solves for disjoint paths.  Post-processing smooths linked
boxes and interpolates internal gaps, but retains the distinction between observed
and synthesized detections.  The pose stage then converts every dense identity track
into a frame-indexed virtual-camera signal.

The components have explicit contracts.  Detection proposes observations; local
association should prioritize fragment purity; global association should prefer a
split over implausible identity evidence; and stabilization must never change track
identity.  This modularity is important because only the association contract is
evaluated against alternative trackers, while the final view also depends on
detector and pose-keypoint quality.
